%=====================================
%   20230301-rits.tex
%   層の超局所理論について
%   toshi2019
%=====================================
\documentclass[dvipdfmx,10pt,aspectratio=169,leqno]{beamer}% dvipdfmxしたい
\title{層の超局所理論について
\subtitle{立命館大学}}
\date{2025年3月1日}
\author{大柴 寿浩}

\institute[NITech]
{
北海道大学大学院理学院数学専攻
}


\usetheme{metropolis}
%\usetheme{Singapore}
\usecolortheme{rose}
\usefonttheme{professionalfonts}
\setbeamertemplate{navigaton symbols}{\false}
\setbeamertemplate{footline}[frame number]
%\usepackage{graphicx,xcolor}





\usepackage{bxdpx-beamer}% dvipdfmxなので必要
\usepackage{pxjahyper}% 日本語で'しおり'したい
\usepackage{minijs}% min10ヤダ
\renewcommand{\kanjifamilydefault}{\gtdefault}
\renewcommand{\emph}[1]{{\upshape\bfseries #1}}

\usepackage{amsmath}
\usepackage{amsthm}
\usepackage{tikz}
\usepackage{color}
\usepackage{ascmac}
\usepackage{amsfonts}
\usepackage{mathrsfs}
\usepackage[mathscr]{eucal}
\usepackage{mathtools}
\usepackage{amssymb}
\usepackage{graphicx}
\usepackage{fancybox}
\usepackage{enumerate}
\usepackage{verbatim}
\usepackage{subfigure}
\usepackage{proof}
\usepackage{listings}
\usepackage{otf}
\usepackage[all]{xy}
\usepackage{amscd}
%\usepackage[dvipdfmx]{hyperref}

\usepackage{xcolor}
\definecolor{darkgreen}{rgb}{0,0.45,0} 
\definecolor{darkred}{rgb}{0.75,0,0}
\definecolor{darkblue}{rgb}{0,0,0.6} 
\hypersetup{
    colorlinks=true,
    citecolor=darkgreen,
    linkcolor=darkblue,
    urlcolor=darkblue,
}

\usetikzlibrary{positioning}


\usepackage{latexsym}
\usepackage{wrapfig}
\usepackage{layout}
\usepackage{url}

\usepackage{okumacro}

\usepackage{comment}

%\usepackage{jdash}
%\usepackage{pxjahyper}
% ----------------------------
% commmand
% ----------------------------
% 執筆に便利なコマンド集です. 
% コマンドを追加する場合は下のスペースへ. 

% 集合の記号 (黒板文字)
\newcommand{\NN}{\mathbb{N}}
\newcommand{\ZZ}{\mathbb{Z}}
\newcommand{\QQ}{\mathbb{Q}}
\newcommand{\RR}{\mathbb{R}}
\newcommand{\CC}{\mathbb{C}}
\newcommand{\PP}{\mathbb{P}}
\newcommand{\KK}{\mathbb{K}}


% 集合の記号 (太文字)
\newcommand{\nn}{\mathbf{N}}
\newcommand{\zz}{\mathbf{Z}}
\newcommand{\qq}{\mathbf{Q}}
\newcommand{\rr}{\mathbf{R}}
\newcommand{\cc}{\mathbf{C}}
\newcommand{\pp}{\mathbf{P}}
\newcommand{\kk}{\mathbf{k}}

% 特殊な写像の記号
\newcommand{\ev}{\mathop{\mathrm{ev}}\nolimits} % 値写像
\newcommand{\pr}{\mathop{\mathrm{pr}}\nolimits} % 射影

% スクリプト体にするコマンド
%   例えば　{\mcal C} のように用いる
\newcommand{\mcal}{\mathcal}

% 花文字にするコマンド 
%   例えば　{\h C} のように用いる
\newcommand{\h}{\mathscr}

% ヒルベルト空間などの記号
\newcommand{\F}{\mcal{F}}
\newcommand{\X}{\mcal{X}}
\newcommand{\Y}{\mcal{Y}}
\newcommand{\Hil}{\mcal{H}}
\newcommand{\RKHS}{\Hil_{k}}
\newcommand{\Loss}{\mcal{L}_{D}}
\newcommand{\MLsp}{(\X, \Y, D, \Hil, \Loss)}

% 偏微分作用素の記号
\newcommand{\p}{\partial}

% 角カッコの記号 (内積は下にマクロがあります)
\newcommand{\lan}{\langle}
\newcommand{\ran}{\rangle}



% 圏の記号など
\newcommand{\Set}{{\bf Set}}
\newcommand{\Vect}{{\bf Vect}}
\newcommand{\FDVect}{{\bf FDVect}}
%\newcommand{\Ring}{{\bf Ring}}
\newcommand{\Ab}{{\mathsf{Ab}}}
\newcommand{\Mod}{\mathop{\mathrm{Mod}}\nolimits}
\newcommand{\Modf}{\mathop{\mathrm{Mod}^\mathrm{f}}\nolimits}
\newcommand{\CGA}{{\bf CGA}}
\newcommand{\GVect}{{\bf GVect}}
\newcommand{\Lie}{{\bf Lie}}
\newcommand{\dLie}{{\bf Liec}}



% 射の集合など
\newcommand{\Map}{\mathop{\mathrm{Map}}\nolimits} % 写像の集合
\newcommand{\Hom}{\mathop{\mathrm{Hom}}\nolimits} % 射集合
\newcommand{\End}{\mathop{\mathrm{End}}\nolimits} % 自己準同型の集合
\newcommand{\Aut}{\mathop{\mathrm{Aut}}\nolimits} % 自己同型の集合
\newcommand{\Mor}{\mathop{\mathrm{Mor}}\nolimits} % 射集合
\newcommand{\Ker}{\mathop{\mathrm{Ker}}\nolimits} % 核
\newcommand{\Img}{\mathop{\mathrm{Im}}\nolimits} % 像
\newcommand{\Cok}{\mathop{\mathrm{Coker}}\nolimits} % 余核
\newcommand{\Cim}{\mathop{\mathrm{Coim}}\nolimits} % 余像

% その他便利なコマンド
\newcommand{\dip}{\displaystyle} % 本文中で数式モード
\newcommand{\e}{\varepsilon} % イプシロン
\newcommand{\dl}{\delta} % デルタ
\newcommand{\pphi}{\varphi} % ファイ
\newcommand{\ti}{\tilde} % チルダ
\newcommand{\pal}{\parallel} % 平行
\newcommand{\op}{{\rm op}} % 双対を取る記号
\newcommand{\lcm}{\mathop{\mathrm{lcm}}\nolimits} % 最小公倍数の記号
\newcommand{\Probsp}{(\Omega, \F, \P)} 
\newcommand{\argmax}{\mathop{\rm arg~max}\limits}
\newcommand{\argmin}{\mathop{\rm arg~min}\limits}





% ================================
% コマンドを追加する場合のスペース 
%\newcommand{\OO}{\mcal{O}}



\makeatletter
\renewenvironment{proof}[1][\proofname]{\par
  \pushQED{\qed}%
  \normalfont \topsep6\p@\@plus6\p@\relax
  \trivlist
  \item[\hskip\labelsep
%        \itshape
         \bfseries
%    #1\@addpunct{.}]\ignorespaces
    {#1}]\ignorespaces
}{%
  \popQED\endtrivlist\@endpefalse
}
\makeatother

\renewcommand{\proofname}{\textrm{Proof.}}



%\renewcommand\proofname{\bf 証明} % 証明
\numberwithin{equation}{subsection}
\newcommand{\cTop}{\textsf{Top}}
%\newcommand{\cOpen}{\textsf{Open}}
\newcommand{\Op}{\mathop{\textsf{Op}}\nolimits}
\newcommand{\Ob}{\mathop{\textrm{Ob}}\nolimits}
\newcommand{\id}{\mathop{\mathrm{id}}\nolimits}
\newcommand{\pt}{\mathop{\mathrm{pt}}\nolimits}
\newcommand{\res}{\mathop{\rho}\nolimits}
\newcommand{\A}{\mcal{A}}
\newcommand{\B}{\mcal{B}}
\newcommand{\C}{\mcal{C}}
\newcommand{\D}{\mcal{D}}
\newcommand{\E}{\mcal{E}}
\newcommand{\G}{\mcal{G}}
%\newcommand{\H}{\mcal{H}}
\newcommand{\I}{\mcal{I}}
\newcommand{\J}{\mcal{J}}
\newcommand{\OO}{\mcal{O}}
\newcommand{\Ring}{\mathop{\textsf{Ring}}\nolimits}
\newcommand{\cAb}{\mathop{\textsf{Ab}}\nolimits}
%\newcommand{\Ker}{\mathop{\mathrm{Ker}}\nolimits}
\newcommand{\im}{\mathop{\mathrm{Im}}\nolimits}
\newcommand{\Coker}{\mathop{\mathrm{Coker}}\nolimits}
\newcommand{\Coim}{\mathop{\mathrm{Coim}}\nolimits}
\newcommand{\rank}{\mathop{\mathrm{rank}}\nolimits}
\newcommand{\Ht}{\mathop{\mathrm{Ht}}\nolimits}
\newcommand{\supp}{\mathop{\mathrm{supp}}\nolimits}
\newcommand{\colim}{\mathop{\mathrm{colim}}}
\newcommand{\Tor}{\mathop{\mathrm{Tor}}\nolimits}

\newcommand{\cat}{\mathscr{C}}

%筆記体
\newcommand{\cA}{\mcal{A}}
\newcommand{\cB}{\mcal{B}}
\newcommand{\cC}{\mcal{C}}
\newcommand{\cD}{\mcal{D}}
\newcommand{\cE}{\mcal{E}}
\newcommand{\cF}{\mcal{F}}
\newcommand{\cG}{\mcal{G}}
\newcommand{\cH}{\mcal{H}}
\newcommand{\cI}{\mcal{I}}
\newcommand{\cJ}{\mcal{J}}
\newcommand{\cK}{\mcal{K}}
\newcommand{\cL}{\mcal{L}}
\newcommand{\cM}{\mcal{M}}
\newcommand{\cN}{\mcal{N}}
\newcommand{\cO}{\mcal{O}}
\newcommand{\cP}{\mcal{P}}
\newcommand{\cQ}{\mcal{Q}}
\newcommand{\cR}{\mcal{R}}
\newcommand{\cS}{\mcal{S}}
\newcommand{\cT}{\mcal{T}}
\newcommand{\cU}{\mcal{U}}
\newcommand{\cV}{\mcal{V}}
\newcommand{\cW}{\mcal{W}}
\newcommand{\cX}{\mcal{X}}
\newcommand{\cY}{\mcal{Y}}
\newcommand{\cZ}{\mcal{Z}}


\newcommand{\scA}{\mathscr{A}}
\newcommand{\scB}{\mathscr{B}}
\newcommand{\scC}{\mathscr{C}}
\newcommand{\scD}{\mathscr{D}}
\newcommand{\scE}{\mathscr{E}}
\newcommand{\scF}{\mathscr{F}}
\newcommand{\scN}{\mathscr{N}}
\newcommand{\scO}{\mathscr{O}}
\newcommand{\scV}{\mathscr{V}}
\newcommand{\scU}{\mathscr{U}}


\newcommand{\ibA}{\mathop{\text{\textit{\textbf{A}}}}}
\newcommand{\ibB}{\mathop{\text{\textit{\textbf{B}}}}}
\newcommand{\ibC}{\mathop{\text{\textit{\textbf{C}}}}}
\newcommand{\ibD}{\mathop{\text{\textit{\textbf{D}}}}}
\newcommand{\ibE}{\mathop{\text{\textit{\textbf{E}}}}}
\newcommand{\ibF}{\mathop{\text{\textit{\textbf{F}}}}}
\newcommand{\ibG}{\mathop{\text{\textit{\textbf{G}}}}}
\newcommand{\ibH}{\mathop{\text{\textit{\textbf{H}}}}}
\newcommand{\ibI}{\mathop{\text{\textit{\textbf{I}}}}}
\newcommand{\ibJ}{\mathop{\text{\textit{\textbf{J}}}}}
\newcommand{\ibK}{\mathop{\text{\textit{\textbf{K}}}}}
\newcommand{\ibL}{\mathop{\text{\textit{\textbf{L}}}}}
\newcommand{\ibM}{\mathop{\text{\textit{\textbf{M}}}}}
\newcommand{\ibN}{\mathop{\text{\textit{\textbf{N}}}}}
\newcommand{\ibO}{\mathop{\text{\textit{\textbf{O}}}}}
\newcommand{\ibP}{\mathop{\text{\textit{\textbf{P}}}}}
\newcommand{\ibQ}{\mathop{\text{\textit{\textbf{Q}}}}}
\newcommand{\ibR}{\mathop{\text{\textit{\textbf{R}}}}}
\newcommand{\ibS}{\mathop{\text{\textit{\textbf{S}}}}}
\newcommand{\ibT}{\mathop{\text{\textit{\textbf{T}}}}}
\newcommand{\ibU}{\mathop{\text{\textit{\textbf{U}}}}}
\newcommand{\ibV}{\mathop{\text{\textit{\textbf{V}}}}}
\newcommand{\ibW}{\mathop{\text{\textit{\textbf{W}}}}}
\newcommand{\ibX}{\mathop{\text{\textit{\textbf{X}}}}}
\newcommand{\ibY}{\mathop{\text{\textit{\textbf{Y}}}}}
\newcommand{\ibZ}{\mathop{\text{\textit{\textbf{Z}}}}}

\newcommand{\ibx}{\mathop{\text{\textit{\textbf{x}}}}}

%\newcommand{\Comp}{\mathop{\mathrm{C}}\nolimits}
%\newcommand{\Komp}{\mathop{\mathrm{K}}\nolimits}
%\newcommand{\Domp}{\mathop{\mathsf{D}}\nolimits}%複体のホモトピー圏
%\newcommand{\Comp}{\mathrm{C}}
%\newcommand{\Komp}{\mathrm{K}}
%\newcommand{\Domp}{\mathsf{D}}%複体のホモトピー圏

\newcommand{\Comp}{\mathsf{C}}
\newcommand{\Komp}{\mathsf{K}}
\newcommand{\Domp}{\mathsf{D}}
\newcommand{\Kompl}{\mathop{\mathsf{K}^\mathrm{+}}\nolimits}
\newcommand{\Kompu}{\mathop{\mathsf{K}^\mathrm{-}}\nolimits}
\newcommand{\Kompb}{\mathop{\mathsf{K}^\mathrm{b}}\nolimits}
\newcommand{\Dompl}{\mathop{\mathsf{D}^\mathrm{+}}\nolimits}
\newcommand{\Dompu}{\mathop{\mathsf{D}^\mathrm{-}}\nolimits}
\newcommand{\Dompb}{\mathop{\mathsf{D}^\mathrm{b}}\nolimits}
\newcommand{\Dompbf}{\mathop{\mathsf{D}_\mathrm{f}^\mathrm{b}}\nolimits}




\newcommand{\CCat}{\Comp(\cat)}
\newcommand{\KCat}{\Komp(\cat)}
\newcommand{\DCat}{\Domp(\cat)}%圏Cの複体のホモトピー圏
\newcommand{\HOM}{\mathop{\mathscr{H}\hspace{-2pt}om}\nolimits}%内部Hom
\newcommand{\RHOM}{\mathop{\mathrm{R}\hspace{-1.5pt}\HOM}\nolimits}

\newcommand{\muS}{\mathop{\mathrm{SS}}\nolimits}
\newcommand{\RG}{\mathop{\mathrm{R}\hspace{-0pt}\Gamma}\nolimits}
\newcommand{\RHom}{\mathop{\mathrm{R}\hspace{-1.5pt}\Hom}\nolimits}
\newcommand{\Rder}{\mathrm{R}}

\newcommand{\simar}{\mathrel{\overset{\sim}{\rightarrow}}}%同型右矢印
\newcommand{\simarr}{\mathrel{\overset{\sim}{\longrightarrow}}}%同型右矢印
\newcommand{\simra}{\mathrel{\overset{\sim}{\leftarrow}}}%同型左矢印
\newcommand{\simrra}{\mathrel{\overset{\sim}{\longleftarrow}}}%同型左矢印

\newcommand{\hocolim}{{\mathrm{hocolim}}}
\newcommand{\indlim}[1][]{\mathop{\varinjlim}\limits_{#1}}
\newcommand{\sindlim}[1][]{\smash{\mathop{\varinjlim}\limits_{#1}}\,}
\newcommand{\Pro}{\mathrm{Pro}}
\newcommand{\Ind}{\mathrm{Ind}}
\newcommand{\prolim}[1][]{\mathop{\varprojlim}\limits_{#1}}
\newcommand{\sprolim}[1][]{\smash{\mathop{\varprojlim}\limits_{#1}}\,}

\newcommand{\Sh}{\mathrm{Sh}}
\newcommand{\PSh}{\mathrm{PSh}}

\newcommand{\rmD}{\mathrm{D}}

\newcommand{\Lloc}[1][]{\mathord{\mathcal{L}^1_{\mathrm{loc},{#1}}}}
\newcommand{\ori}{\mathord{\mathrm{or}}}
\newcommand{\Db}{\mathord{\mathscr{D}b}}

\newcommand{\codim}{\mathop{\mathrm{codim}}\nolimits}



\newcommand{\gld}{\mathop{\mathrm{gld}}\nolimits}
\newcommand{\wgld}{\mathop{\mathrm{wgld}}\nolimits}


\newcommand{\tens}[1][]{\mathbin{\otimes_{\raise1.5ex\hbox to-.1em{}{#1}}}}
\newcommand{\etens}{\mathbin{\boxtimes}}
\newcommand{\ltens}[1][]{\mathbin{\overset{\mathrm{L}}\tens}_{#1}}
\newcommand{\mtens}[1][]{\mathbin{\overset{\mathrm{\mu}}\tens}_{#1}}
\newcommand{\lltens}[1][]{{\mathop{\tens}\limits^{\mathrm{L}}_{#1}}}
\newcommand{\letens}{\overset{\mathrm{L}}{\etens}}
\newcommand{\detens}{\underline{\etens}}
\newcommand{\ldetens}{\overset{\mathrm{L}}{\underline{\etens}}}
\newcommand{\dtens}[1][]{{\overset{\mathrm{L}}{\underline{\otimes}}}_{#1}}

\newcommand{\blk}{\mathord{\ \cdot\ }}
\newcommand{\mres}[2][]{{\left.{#1}\right\rvert}_{#2}}


%\newcommand{\hocolim}{{\mathrm{hocolim}}}
%\newcommand{\indlim}[1][]{\mathop{\varinjlim}\limits_{#1}}
%\newcommand{\sindlim}[1][]{\smash{\mathop{\varinjlim}\limits_{#1}}\,}
%\newcommand{\Pro}{\mathrm{Pro}}
%\newcommand{\Ind}{\mathrm{Ind}}
%\newcommand{\prolim}[1][]{\mathop{\varprojlim}\limits_{#1}}
%\newcommand{\sprolim}[1][]{\smash{\mathop{\varprojlim}\limits_{#1}}\,}
\newcommand{\proolim}[1][]{\mathop{\text{\rm``{$\varprojlim$}''}}\limits_{#1}}
\newcommand{\sproolim}[1][]{\smash{\mathop{\rm``{\varprojlim}''}\limits_{#1}}}
\newcommand{\inddlim}[1][]{\mathop{\text{\rm``{$\varinjlim$}''}}\limits_{#1}}
\newcommand{\sinddlim}[1][]{\smash{\mathop{\text{\rm``{$\varinjlim$}''}}\limits_{#1}}\,}
\newcommand{\ooplus}{\mathop{\text{\rm``{$\oplus$}''}}\limits}
\newcommand{\bbigsqcup}{\mathop{``\bigsqcup''}\limits}
\newcommand{\bsqcup}{\mathop{``\sqcup''}\limits}
\newcommand{\dsum}[1][]{\mathbin{\oplus_{#1}}}

\newcommand{\Fct}{\mathop{\mathsf{Fct}}\nolimits}
\newcommand{\Red}{\mathop{\mathsf{Red}}\nolimits}
\newcommand{\Cone}{\mathop{\mathsf{Cone}}\nolimits}
\newcommand{\epi}{\mathop{\mathsf{epi}}\nolimits}




\newcommand{\mtan}[1][]{T\hspace{-1.75pt}{#1}}
\newcommand{\mtp}[1][]{{}^{t\!}{#1}} % transpose of the morphism
%\newcommand{\mpb}[1][]{{}^{t\!}{#1}'} % pull-back of the morphism
\newcommand{\mpb}[1][]{{#1}_{\pi}} % pull-back of the morphism
\newcommand{\mdiff}[1][]{{#1}_{d}} % pull-back of the morphism
\newcommand{\mptan}[1][]{{#1}_{\tau}} % pull-back of the morphism


%\newcommand{\pb}[1][]{\mathbin{\mathop{\times}\limits_{#1}}}

\newcommand{\cotb}[1][]{T^{\ast}{#1}}
\newcommand{\conm}[2][]{T_{#1}^{\hspace{0.7pt}\ast}{#2}}
\newcommand{\norb}[2][]{T_{#1}{#2}}

\newcommand{\cotz}[1][]{T^{\ast}_{#1}{#1}}

\newcommand{\cotn}[1][]{\dot{T}^{\ast}{#1}}


\newcommand{\muhom}{\mu hom}
\newcommand{\muHom}[1][]{\mathrm{Hom}^\mu_{\raise1.5ex\hbox to.1em{}#1}}

\newcommand{\mucat}[2][]{\mathsf{D}^{\mathrm{b}}_{#1}(#2)}
%\newcommand{\mucat}[2][]{\mathsf{D}^{\mathrm{b}}_{#1}(#2)}

\newcommand{\conv}[1][]{\mathbin{\mathop{\circ}\limits_{#1}}}

\newcommand{\pb}[1][]{\mathbin{\mathop{\times}\limits_{#1}}}


%================================================
% 自前の定理環境
%   https://mathlandscape.com/latex-amsthm/
% を参考にした
\newtheoremstyle{mystyle}%   % スタイル名
    {5pt}%                   % 上部スペース
    {5pt}%                   % 下部スペース
    {}%              % 本文フォント
    {}%                  % 1行目のインデント量
    {\bfseries}%                      % 見出しフォント
    {.}%                     % 見出し後の句読点
    {12pt}%                     % 見出し後のスペース
    {\thmname{#1}\thmnumber{ #2}\thmnote{{\hspace{2pt}\normalfont (#3)}}}% % 見出しの書式

\theoremstyle{mystyle}
\newtheorem{AXM}{公理}%[section]
\newtheorem{DFN}[AXM]{定義}
\newtheorem{THM}[AXM]{定理}
\newtheorem*{THM*}{定理}
\newtheorem{PRP}[AXM]{命題}
\newtheorem{LMM}[AXM]{補題}
\newtheorem{CRL}[AXM]{系}
\newtheorem{EG}[AXM]{例}
\newtheorem*{EG*}{例}
\newtheorem{RMK}[AXM]{注意}
\newtheorem{CNV}[AXM]{約束}
\newtheorem{CMT}[AXM]{コメント}
\newtheorem*{CMT*}{コメント}
\newtheorem{NTN}[AXM]{記号}
\newtheorem{CLM}[AXM]{Claim}
\newtheorem{Prop}[AXM]{Proposition}

% 定理環境ここまで
%====================================================

\usepackage{framed}
\definecolor{lightgray}{rgb}{0.75,0.75,0.75}
\renewenvironment{leftbar}{%
  \def\FrameCommand{\textcolor{lightgray}{\vrule width 4pt} \hspace{10pt}}% 
  \MakeFramed {\advance\hsize-\width \FrameRestore}}%
{\endMakeFramed}
\newenvironment{redleftbar}{%
  \def\FrameCommand{\textcolor{lightgray}{\vrule width 1pt} \hspace{10pt}}% 
  \MakeFramed {\advance\hsize-\width \FrameRestore}}%
 {\endMakeFramed}



\renewcommand{\Re}{\mathop{\mathrm{Re}}\nolimits}


% =================================





% ---------------------------
% new definition macro
% ---------------------------
% 便利なマクロ集です

% 内積のマクロ
%   例えば \inner<\pphi | \psi> のように用いる
\def\inner<#1>{\langle #1 \rangle}

% ================================
% マクロを追加する場合のスペース 

%=================================







\def\fracinline#1/#2{\mbox{\raise0.5ex\hbox{\footnotesize$#1$}{\hskip-.1em$/$\hskip-.1em}\raise-0.5ex\hbox{\footnotesize$#2$}}}

\begin{document}




\begin{frame}
    \titlepage
\end{frame}
\begin{comment}
    \begin{frame}\frametitle{書くこと}
    \begin{itemize}
        \item 特に興味を持った定理や理論の背景
        \item それを記述するための記号や概念の導入
        \item 正確な主張の紹介
        \item 証明の基本的なアイデアやアウトラインの紹介
        \item または定理の過程を満たす具体例や定理の応用例の紹介
    \end{itemize}
    \end{frame}
\end{comment}
%\setcounter{section}{3}
%\section*{はじめに}
\begin{frame}
    \frametitle{発表内容}
    \tableofcontents
\end{frame}

\section[層]{層理論}
\subsection{定義}

\begin{frame}
    \frametitle{設定と記号}
    \begin{itemize}
        \item \(X\)：位相空間
        \item \(\Op_X\)：\(X\)の開集合からなる順序集合
        \item \(I_P\)：\(P\in X\)の開近傍からなる有向順序集合
        \item \(U\in \Op_{X}\)：開集合
        \item \(C^0(U)\)：\(U\)上の実数値連続関数環
        \item \((U_i)_{i\in I}\)：\(U\)の被覆．
        \item \(\kk\)：単位元をもつ結合的な環（コホモロジー次元は有限とする）
    \end{itemize}
\end{frame}

\begin{frame}
    \frametitle{簡単な例}

    \textbf{連続関数に対する2つの操作}：
    \begin{description}
        \item[制限]\(f\in C^0(U)\)の\(V\subset U\)上の\(f\rvert_V\in C^0(V)\)への制限．
        \[
            \text{\(U\supset V\supset W\)ならば，}
            \left.({f\rvert_V})\right\rvert_W=f\rvert_W.
        \]
        \item[はりあわせ]
        \((f_i)_{i}\in\prod_{i\in I}^{}C^{0}(U_i)\)で\(f_i=f_j\) on \(U_i\cap U_j\)のとき，
        \(f\rvert_{U_i}=f_i\)として\(f\in C^{0}(U)\)に貼り合わせ．
    \end{description}

    \bigskip
    \begin{itemize}
        \item 制限\(\to\)前層
        \item 貼り合わせ\(\to\)層
    \end{itemize}
    として定式化
\end{frame}



\begin{frame}
    \frametitle{前層の定義}

    \begin{definition}[前層]
        \(\kk\)加群の前層\(F\)とは\(\kk\)加群の族と\(\kk\)加群の射の族の組\[\left(
        \left(F(U)\right)_{U\in\Op(X)},
        \left(\rho_{VU}\colon F(U)\to F(V)\right)_{(V\hookrightarrow U)}
    \right)\]で次をみたすもののこと
    \begin{enumerate}
        %\renewcommand{\labelenumi}{({\arabic{enumi}})}
        \item $\rho_{UU} = \id_{F(U)}$ 
        \item $W \subset V \subset U$ ならば, 
        $\rho_{WU} = \rho_{WV} \circ \rho_{VU}$.\label{enum:res}
    \end{enumerate}
    \end{definition}
\end{frame}
\begin{frame}
    \frametitle{前層の定義}

    \begin{definition}[前層の射]
        \(\kk\)加群の前層の射\(
            \phi\colon F\to G
        \)とは，
        \(\kk\)加群の射の族\[
            (\phi_{U}\colon F(U)
            \to G(U))_{U\in\Op(X)}
        \]
        で次を可換にするもののこと
        \[\vcenter{\xymatrix@C=26pt@R=26pt{
            F(U)
            \ar[r]^-{\phi_{U}}
            \ar[d]^-{\rho^{F}_{VU}}
            &
            {G}(U)
            \ar[d]^-{\rho^{G}_{VU}} 
            \\
            F(V)
            \ar[r]^-{\phi_{V}}
            &
            G(V)
            }}
        \]
    \end{definition}
\end{frame}
\begin{frame}
    \frametitle{前層の例}
\begin{exampleblock}{例}
    \(X\)はそれぞれの例に応じたよい位相空間とする．
    \begin{enumerate}
        \item \(C_X^r\colon U\mapsto C_X^r(U)\coloneqq\left\{\text{\(U\)上の\(C^r\)級関数}\right\}\)は\(X\)上の
        前層である．
        \item \(
            \mathscr{O}_X\colon U\mapsto 
            \mathscr{O}_X(U)\coloneqq\left\{\text{\(U\)上の正則関数}\right\}
        \)は\(X\)上の前層である．
        \item \(
            \mathscr{M}_X\colon U\mapsto 
            \mathscr{M}_X(U)\coloneqq\left\{\text{\(U\)上の有理型関数}\right\}
        \)は\(X\)上の前層である．
        \item \(M\)をアーベル群とする．
        \(U\mapsto M\)は\(X\)上の前層である．
    \end{enumerate}
\end{exampleblock}
\end{frame}

\begin{frame}
    \frametitle{前層の例}
    \begin{exampleblock}{ベクトル束の切断}
        \(E\to X\)を\(C^{\infty}\)ベクトル束とする．
        \[
            U\mapsto \Gamma(U;E)\coloneqq
            \left\{s\colon U\to E;s\colon C^{\infty},\pi\circ s=\id_U\right\}
        \]は前層である．
    \end{exampleblock}
\end{frame}

\begin{frame}
    \frametitle{層の定義}
    \begin{definition}[層]\label{def:sheaf}
        \(F\)：\(X\)上の\(\kk\)加群の前層．
        
        \(F\)が\textbf{層} (sheaf) であるとは，
        \(\forall U\in\Op(X)\)と
        その開被覆\(\left(U_i\right)_{i\in I}\)に対して，
        次の条件が成り立つこと．
        \begin{enumerate}[(S1)]
            \setcounter{enumi}{-1}
            \item \(F(\varnothing)=0\).
            \item \(s\in F(U)\)が各\(i\in I\)に対して\(U_i\)上\(s\rvert_{U_{i}}=0\)ならば，
            \(s=0\).\label{sheaf-cond1}
            \item \(\left(s_i\right)_i\in\prod_{i}F(U_i)\)が
            各\(i,j\in I\)で\(U_i\cap U_j\ne\varnothing\)となるものに対して\(U_i\cap U_j\)上\(s_i\rvert_{U_i\cap U_j}-s_j\rvert_{U_i\cap U_j}=0\)ならば，
            \(s\in F(U)\)で各\(U_i\)上\(s\rvert_i=s_i\)となるものが存在する．\label{sheaf-cond2}
        \end{enumerate}
    \end{definition}    
\end{frame}

\begin{frame}
    \frametitle{層の定義}
    層の射は前層としての射

    \bigskip
    記号を定める．
    \begin{itemize}
        \item \(\PSh(X)\coloneqq\left\{\text{\(X\)上の前層}\right\}\)
        \item \(\Sh(X)\coloneqq\left\{\text{\(X\)上の層}\right\}\)
        \item \(\Mod(\kk_X)\coloneqq\left\{\text{\(X\)上の\(\kk\)加群の層}\right\}\)
        \item \(\Ab\coloneqq\left\{\text{アーベル群}\right\}\)
    \end{itemize}
\end{frame}

\begin{frame}
    \frametitle{切断関手}

    \begin{DFN}[切断関手]
        %\(F\in\Sh(X)\), 
        \(U\in\Op(X)\)：fixed
        \[
            \Gamma(U;\blk)\colon \Sh(X)\to \Ab;\quad F\mapsto
            \Gamma(U;F)=F(U)
        \]
    \end{DFN}    
\end{frame}



\begin{frame}
    \frametitle{層の例}
    \begin{EG}\label{eg:sheaf}
        \(X\)を位相空間とする．
        \begin{enumerate}
            \item 前層\(C_X^r\)は\(X\)上の層である．
            \item 前層\(\mathscr{O}_X\)は\(X\)上の層である．
            \item 前層\(\mathscr{M}_X\)は\(X\)上の層である．
            \item \(M\)をアーベル群とする．前層\(M\colon U\mapsto M\)は\(X\)上の層ではない．
            実際，\(X\)を2点からなる離散空間\(2=\left\{0,1\right\}\)とし，その上の
            前層として\(M=\cc\)を考えると，
            \(1\in\cc(\{0\})\), \(i\in\cc(\{1\})\)に対して，
            \(\cc(\{0\}\cap\{1\})=\cc(\varnothing)=0\)
            より，\(\rho_{\varnothing,{0}}(1)=0
            =\rho_{\varnothing,{1}}(i)\)となるが，
            \(2\)上の切断\(z\in\cc(2)\)で\(z\rvert_{\{0\}}=1,
            z\rvert_{\{1\}}=i\)をみたすものは存在しない．
            したがって，条件 (S\ref{sheaf-cond2}) が成り立たない．\label{eg:non-sheaf}
        \end{enumerate}
    \end{EG}    
\end{frame}

\begin{frame}
    \frametitle{前層の例}
    \begin{EG}
        \(E\to X\)を\(C^{\infty}\)ベクトル束とする．
        \[
            C^{\infty}(E)\colon U\mapsto \Gamma(U;E)\coloneqq
            \left\{s\colon U\to E;s\colon C^{\infty}, \pi\circ s=\id_U\right\}
        \]は層である．
    \end{EG}
\end{frame}


\subsection{層化と茎}

\begin{frame}
    \frametitle{層化}
    層ではない前層に対して，
    「最も近い」層を自然に対応させる．
    \begin{THM}\label{thm:sheafification}
        \(X\)上の任意の前層\(F\)に対して，
        \(X\)上の層\(F^\dagger\)と
        前層の射\(\theta\colon F\to F^\dagger\)で
        次の普遍性をみたすものが
        ただ一つ存在する．
        \begin{quote}
            任意の層 \({G}\in\Sh(X)\)と
            前層の射\(\varphi\colon F\to{G}\)に
            対し，層の射\(\varphi^\dagger\colon F^\dagger
            \to{G}\)で，
            \[\varphi^\dagger\circ \theta=\varphi\]をみたすものが
            ただ一つ存在する．
        \end{quote}
    \end{THM}    
\end{frame}

\begin{frame}
    \frametitle{茎の例から}
    \begin{itemize}
        \item \(X\)：例えばリーマン面，\(U\subset X\)：開集合
        \item 正則関数\(f\in\OO_{X}(U)\)は
        各点\(P\)のまわりで冪級数展開可能
        \item 関数\(f\in \mathscr{O}_X(U)\)と\(g\in \mathscr{O}_X(V)\)が\(P\)の近くで等しいとは
        
        \begin{quote}
            \(U\cap V\)に含まれる\(P\)の近傍\(W\)とって，
            \(P\)のまわりの座標\(z\colon W\to z(W)\)
            を用いて\(f,g\)を展開したとき，
            収束冪級数として等しいということ
        \end{quote}
    \end{itemize}

定義域の異なる関数（環）に対して，
各点まわりに注目するという操作を茎として定式化．

\end{frame}

\begin{frame}
    \frametitle{茎の定義}
    \begin{DFN}[茎]
        \begin{itemize}
            \item \(F\in\Sh(X)\)
            \item \(P\in X\)
        \end{itemize}
        \(F\)の\(P\)での\textbf{茎} (stalk) \(
            F_P
        \)を次で定める．
        \[
            F_P\coloneqq
            \varinjlim_{U\in I_P}F(U).
        \]
    
        \(F(U)\to F_P\)から定まる
        \(s\in F(U)\)の行き先を\(s_P\)とかき，
        \(s\)の\textbf{芽} (germ) という．
    \end{DFN}
\end{frame}

\begin{frame}
    \frametitle{茎の別の定義}
    \[
        F_{P}\cong
        \left.\left(
            \bigsqcup_{U\in I_P}F(U)
        \right)\right/\mathord{\sim}
    \]
ここで，\((f,U)\), \((g,V)\in \bigsqcup_{U\in I_P}F(U)\)に
対し，\[
    (f,U)\sim(g,V)
    \colon\Longleftrightarrow
    \begin{cases}
        \text{\(P\)の開近傍\(W\in I_P\)で\(W\subset U\cap V\)と}\\
        \text{\(f\rvert_W=g\rvert_W\)をみたすものが存在する．}
    \end{cases}
\]
つまり，\(P\)のまわりでの挙動が同じ切断を
同一視した類が芽．
\end{frame}


\begin{frame}
    \frametitle{層化の構成}

    開集合\(U\)に対し，
    \[
        F^\dagger(U)\coloneqq
        \left\{
            (s^P)_{P\in U}\in\prod_{P\in U}^{}F_P;
            \begin{aligned}
                \text{任意の\(P\in U\)に対し，近傍\(W\in I_P\cap U\)と}\\
                \text{切断\(t\in F(W)\)で，任意の\(Q\in W\)に対し，}\\            
                \text{\(s^Q=t_Q\)となるものが存在する．}            
            \end{aligned}
        \right\}
    \]
    とおき，
    制限射\(F^\dagger(U)\to F^\dagger(V)\)を
    \((s^P)_{P\in U}\mapsto(s^P)_{P\in V}\)で定めると，
    前層\(F^\dagger\)が定まり，
    \(F^\dagger\)が層の条件をみたすことも確かめられる．

    \begin{CMT}層化は茎を保つ．
        \(F_P={F}^{\dagger}_{P}\)
    \end{CMT}
    
\end{frame}

\begin{frame}
    \frametitle{層化の例}

    \begin{exampleblock}{定数層}
        \(M\colon U\mapsto M\)をアーベル群\(M\)から定まる前層とする．
        \(M\)の層化\(M^\dagger\)は\[
            M^\dagger(U)=\left\{
            \text{\(U\)上の\(M\)に値をとる局所定数関数}\right\}
        \]である．
        \(M^\dagger\)を\(M_X\)とかき，\(X\)上の定数層と呼ぶ．
    \end{exampleblock}

    \begin{PRP}
        \[M_X(U)\cong M^{\# \pi_0(U)}\]
        ここで，\(\pi_0(U)\)は\(U\)の連結成分の集合．
    \end{PRP}

\end{frame}

\begin{frame}
    \frametitle{層の例}
    \begin{exampleblock}{定数層の\(0\)による延長}\label{eg:const}
        \(X\)：位相空間，\(Z\subset X\)：閉集合
        
        \(\kk\)：体
        \[
            U\mapsto
            \kk_{X,Z}(U)\coloneqq\left\{
                f\colon Z\cap U\to\kk; 
                \text{\(f\)は局所定数関数}
            \right\}
        \]
    \end{exampleblock}
    %一般にベクトル束は局所自由層と対応する．
\end{frame}
\begin{frame}
    \frametitle{\(Z\)上の定数層の\(0\)による延長のイメージ}
\end{frame}

\section[層係数コホモロジー]{層係数コホモロジー}

\begin{frame}
    \frametitle{層係数コホモロジー：前置き}
    空間\(X\)に対し，そのコホモロジー
    \[
        H^i(X;G)\quad\text{（\(G\)はアーベル群）}
    \]は重要な量．
    
    場所によって係数\(G\)を取り替えると，
    局所系や層を係数とするコホモロジー
    \[
        H^i(X;F)
    \]が出てくる．
    （\(
        \text{局所系}
        \leftrightarrow
        \text{局所定数層}
        \subset\text{層}
    \)）
\end{frame}

\begin{frame}
    \frametitle{層係数コホモロジー}
    コホモロジーを計算するときには，``良い層''\(\mathscr{I}^i\)たちで分解して
    \[\vcenter{\xymatrix@C=22pt@R=26pt{
        \text{もとの層}F
        \ar@{~>}[d]^-{\text{分解}}
        \\ 
        \text{層の複体\(\mathscr{I}^\bullet\)}
        \ar@{~>}[d]^-{H^i}
        \\        
        \text{コホモロジー}H^i(X;F)
    }}\]
    のようにする
\end{frame}

\begin{frame}
    \frametitle{層係数コホモロジー}
    例えば，切断関手\(\Gamma(U;\blk)\)は左完全：
    \[\vcenter{\xymatrix{
        0\ar[r]& F\ar[r]& G\ar[r]\ar@{=>}[d]& H\ar[r]& 0& \text{exact in \(\Sh(X)\)}
        \\
        0\ar[r]&\Gamma(U;F)\ar[r]&\Gamma(U;G)\ar[r]&\Gamma(U;H)& & \text{exact in \(\Ab\)}
    }}\]
    したがって，その完全でない度合いを測る，
    右導来関手\(\RG(U;\blk)\)が定まる．
    
    この\(i\)次の部分\(\Rder^i\Gamma(U;F)\)などが
    層係数の\(H^i(U;F)\)など．
\end{frame}

\begin{frame}
    \frametitle{層の導来圏}
    導来圏の理論を用いると見通しがよい
    \begin{NTN}
        層の有界な複体を対象とし，
        コホモロジー間の同型を誘導する射（擬同型）を可逆にした
        圏を\(\Dompb(\kk_X)\)で表す．
    \end{NTN}
    \(\Dompb(\cc_X)\)では例えば定数層のド・ラーム分解が同型になる：
    \[\vcenter{\xymatrix{
        0\ar[r]&
        \cc_X \ar[r] \ar[d]^-{\mathrm{qis}}&
        0\\
        0\ar[r]&
        \Omega_X^0\ar[r]&
        \Omega_X^1\ar[r]&
        \cdots\ar[r]&
        \Omega_X^n\ar[r]&0
    }}\]
\end{frame}

\begin{frame}
    \frametitle{完全三角}
    導来圏では完全三角 (distinguished triangle) が完全列の代わりになる．
    \[
        F\to G\to H\to+1\quad\mathrm{d.t.}
    \]
    とかいたとき，長完全列
    \begin{align*}
        0&\to H^0(F)\to H^0(G)\to H^0(H)\\
        &\to H^1(F)\to H^1(G)\to H^1(H)\\
        &\to \cdots
    \end{align*}
    が得られる．つまり，上の完全三角で\(H=0\)なら，
    \[
    F\simarr G\quad \mathrm{qis}
    \]となる．
\end{frame}



\begin{frame}
    \frametitle{層の六則演算}
    \(f\colon M\to N\)を多様体の射とする．
    このとき，次の6つの関手が定まる．
    \begin{align*}
        \Rder{f_\ast}&\colon\Dompb(\kk_M)\to\Dompb(\kk_N)\\
        \Rder{f_!}&\colon\Dompb(\kk_M)\to\Dompb(\kk_N)\\
        f^{-1}&\colon\Dompb(\kk_N)\to\Dompb(\kk_M)\\
        f^{!}&\colon\Dompb(\kk_N)\to\Dompb(\kk_M)\\
        \RHOM&\colon\Dompb(\kk_M)^{\op}\times\Dompb(\kk_M)\to\Dompb(\kk_M)\\
        \lltens[]&\colon\Dompb(\kk_M^{\op})\times\Dompb(\kk_M)\to\Dompb(\kk_M)
    \end{align*}
    \begin{alertblock}{注意}
        \(f^!\)は導来圏の間の関手としてしか定まらない．        
    \end{alertblock}
\end{frame}

\begin{frame}
    \frametitle{六則演算における公式の例}
    \begin{exampleblock}{順像と切断}
        \(f\colon M\to N\)を多様体の射，
        \(B\)を\(N\)の局所閉集合とすると，
        \[
            \RG(B;\Rder{f_\ast}F)\cong\RG(f^{-1}(B);F).
        \]
    \end{exampleblock}
    \begin{itemize}
        \item 左辺は\(F\)に\(\Rder{f_\ast}\)と\(\RG(B;\blk)\)の合成を適用したもの．
        \item 右辺は\(F\)に\(\RG(f^(B);\blk)\)を適用したもの．
        \item \(\RG(X;\blk)\)は一点への写像\(
            a_X\colon X\to\{\pt\}
        \)の順像\({a_{X}}_\ast\)の右導来関手\(\Rder{{a_{X}}_\ast}\)と同じ．
    \end{itemize}
    \begin{alertblock}{コメント}
        導来関手の合成がうまくいくのも導来圏を考えるメリット．
    \end{alertblock}
\end{frame}

\begin{frame}
    \frametitle{局所コホモロジーと\(0\)による延長}
    \begin{DFN}
        局所閉集合\(i\colon Z\hookrightarrow M\)での
        局所コホモロジー\(\RG_{Z}\)と\(0\)による延長\((\blk)_Z\)を
        それぞれ\begin{align*}
            \RG_{Z}&\coloneqq i_{\ast}i^{!}\colon \Dompb(\kk_M)\to\Dompb(\kk_M)\\
            (\blk)_Z&\coloneqq i_{!}i^{-1}\colon \Dompb(\kk_M)\to\Dompb(\kk_M)
        \end{align*}
        で定める．
    \end{DFN}
    %\(f\colon M\to N\)を多様体の射とする．
    %このとき，
    \begin{exampleblock}{切除三角}
        \(Z\subset M\)を閉集合とする．このとき，つぎは完全三角．
        \[
            \RG_{Z}(F)\to F\to \RG_{M-Z}(F)\to+1
        \]
    \end{exampleblock}
\end{frame}


\section{層の超局所台}
\newcommand{\Char}{\mathop{\mathrm{Char}}\nolimits}
\newcommand{\singsupp}{\mathop{\mathrm{singsupp}}\nolimits}
\newcommand{\SingSpec}{\mathop{\mathrm{S.S.}}\nolimits}

\begin{frame}\frametitle{``超局所''?}
    \[
        \text{``超局所''}=\text{余接束上で局所}
    \]
    いったん，古典的な微分方程式の問題で説明する．
    \begin{itemize}
        \item 開集合\(\varOmega\subset \rr^n\)に対し\(\cotn[\varOmega]\coloneqq \cotb[\varOmega]-\cotz[\varOmega]\cong \varOmega\times(\rr^n-\{0\})\to \varOmega\)
        \item ``超関数''\(u\)に対し\(\singsupp(u)\coloneqq\left\{x\in\varOmega;\text{\(u(x)\)は\(x\)で解析的でない}\right\}\)
        \item 微分作用素\(P\)に対し\(\Char(P)\coloneqq\{(x;\xi)\in\cotn[\varOmega];\sigma(P)(x;\xi)=0\}\)（\(\sigma(P)\)は\(P\)の"主要表象"）
        \item \(u\)の特異性スペクトル：\(\SingSpec(u)\subset\cotn[\varOmega]\)は\(
            \pi(\SingSpec(u))=\singsupp(u)
        \)をみたす集合
    \end{itemize}
\end{frame}

\begin{frame}\frametitle{``超局所''?}
    \begin{enumerate}[(i)]
    \item \((x,y)\)平面上のラプラス作用素\[
        \Delta=\partial^2_x+\partial^2_y
    \]
    の主要表象は\(\sigma(\Delta)=\xi_1^2+\xi_2^2\).
    特性多様体は\(\Char(\Delta)=\{\xi_1^2+\xi_2^2\}=\varnothing\).
    \item \((t,x,y)\)空間上の
    拡散作用素\[
        P=\partial_t-\Delta
    \]
    の主要表象は\(\sigma(P)=-\xi_2^2-\xi_3^2\). 
    特性多様体は\(
        \Char(P)=
        \{
            (t,x,y;\xi_1,\xi_2,\xi_3);
            \xi_1^2+\xi_2^2
        \}
        =\varnothing
    \).
    \item \((t,x,y)\)空間上の
    ダランベール作用素\[
        \Box=\partial^2_t-\Delta
    \]
    の主要表象は\(\sigma(\Box)=\xi_1^2-\xi_2^2-\xi_3^2\). 
    特性多様体は\(\Char(\Delta)=\{\xi_1^2+\xi_2^2\}=\varnothing\).
    \end{enumerate}
\end{frame}

\begin{frame}\frametitle{超局所解析における定理} 
    \begin{THM}[佐藤の基本定理]
        \(f\)：既知関数，
        \(P\)：実解析的係数微分作用素とする．
        \(u\)が微分方程式\(Pu=f\)の解のとき
        \[
            \SingSpec(u)\subset \Char(P)\cup\SingSpec(f).
        \]
    \end{THM}
    これを\(f=0\), \(P\)：楕円型の場合に適用すると
    \[
        \Char(P)\cup \SingSpec(f)=\varnothing
    \]
    なので\(\singsupp(u)=\pi(\SingSpec(u))=\varnothing\)．
    すなわち，
    \begin{alertblock}{楕円型正則性定理}
        実解析的係数の楕円型方程式の解は解析的である．
    \end{alertblock}
\end{frame}


\begin{frame}\frametitle{層の超局所台 ------ 定義の動機づけ}
    %\(\Dompb(\kk_X)\ni F\mapsto \SS(F)\subset T^{\ast}X\)
    \(F\in\Dompb(\kk_M)\)の超局所台\(\muS(F)\)を定義したい．

    \(p=(x_0;\xi_0)\in\cotb[M]\)とする．
    \(M\)上の実\(C^1\)級関数\(\varphi\)で\(
        d\varphi(x_0)=\xi_0
    \)となるものを考える．

    \(i\)次コホモロジーにおける``方向\(p\)への制限''
    \[
        H^iF_{x_0}\cong \indlim[U\in I_{x_0}]H^i(U;F)
        \to
        \indlim[U\in I_{x_0}]H^i(U\cap\{x;\varphi(x)<\varphi(x_0)\};F).
    \]
    が同型になるというのが，\(p\)の方向にコホモロジーが伝播するということ．
    三角で考えると，次の定義になる．
\end{frame}


\begin{frame}\frametitle{層の超局所台 ------ 定義}
    %\(\Dompb(\kk_X)\ni F\mapsto \SS(F)\subset T^{\ast}X\)
    
    \begin{Definition}[層の超局所台]
        \((x_0;\xi_0)\notin \SS(F)\)となるのは，
        \(U\in I_{p}\)で，任意の\(x_1\in X\)と
        \(x_1\)の近傍で定義された，\(
            d\varphi(x_1)\in U
        \)となる\(\varphi\)に対し，\[
            \left(\RG_{\{\varphi<\varphi(x_0)\}}(F)\right)_{x_1}=0
        \]
        となるものが存在するときであるとする．
    \end{Definition}
\end{frame}

\begin{frame}
    \frametitle{
        超局所台の例：\(Z\coloneqq\{y\leqq0\}\subset\rr^2\)に
        台を持つ定数層\(\kk_{\rr^2,Z}\)の場合}
\end{frame}

\begin{frame}
    \frametitle{
        超局所台の例：カスプのグラフに台を持つ定数層の場合}
\end{frame}


\begin{frame}\frametitle{``超局所''?}
    次の図式が基本的．
    \[\vcenter{\xymatrix@C=22pt@R=26pt{
        TM
        \ar[d]^-{\tau_M}
        \ar[r]^-{f'}
        &
        M\pb[N] TN
        %\ar@{}[dr]|\square
        \ar[d]_-{}
        \ar[r]^-{\mptan[f]}
        &
        TN
        \ar[d]^-{\tau_N}
        &
        \cotb[M]
        \ar[d]^-{\pi_M}
        &
        M\pb[N] \cotb[N]
        %\ar@{}[dr]|\square
        \ar[d]_-{}
        \ar[r]^-{\mpb[f]}
        \ar[l]_-{\mdiff[f]}
        &
        \cotb[N]
        \ar[d]^-{\pi_N}
        \\
        M
        \ar@{=}[r]
        &
        M
        \ar[r]^-{f}
        %\ar@{}[dr]|\square
        &
        N,
        &
        M
        \ar@{=}[r]
        &
        M
        \ar[r]^-{f}
        %\ar@{}[dr]|\square
        &
        N.
    }}\]
\end{frame}
\begin{frame}
    \frametitle{超局所台の関手的性質}
    \begin{THM}[{\cite[5.4節]{KS90}}参照]\label{thm15}
        \(f\colon M\to N\)を多様体の射とし，
        \(F\in\Dompb(\kk_M)\)，\(G\in\Dompb(\kk_N)\)とする．
        \begin{enumerate}[(i)]
            \item \(f\)が\(\supp(F)\)上固有ならば，
            \(
                \muS(\Rder{f_{!}})
                \subset
                \mpb[f]\mdiff[f]^{-1}\muS(F)
            \)である.\label{item15-1}
            \item \(f\)が\(\muS(G)\)上非特性的ならば，\(
                \muS(f^{-1}G)\subset\mdiff[f]\mpb[f]\muS(G)
            \)である.\label{item15-2}
            \item \(f\)が平滑であるとする．
            このとき，\(
                \muS(F)\subset M\mathbin{\times_{N}} \cotb[N]
            \)となることと，
            すべての\(j\in\zz\)に対し層\(H^j(F)\)が\(f\)の
            ファイバー上で局所定数層となることは同値である．\label{item15-3}
        \end{enumerate}
    \end{THM}
\end{frame}

\section{モース理論}


\begin{frame}
    \frametitle{モース理論の結果}
    \begin{itemize}
        \item \(M_t\coloneqq\left\{x\in M;\psi(x)<t\right\}\)
        \item \(F\in\Dompb(\kk_M)\)
    \end{itemize}
    とする．
    \begin{THM}
        \(\psi\colon M\to \rr\)を\(a\leqq\psi(x)< b\)で
        臨界値をとらない関数とする．
        このとき，\[H^j(M_a;\rr)\cong H^j(M_b;\rr)\]
    \end{THM}
\end{frame}

\begin{frame}
    \frametitle{層係数への拡張}
    \begin{PRP}[超局所モースの補題]
        \(\psi\colon M\to \rr\)を\(a\leqq\psi(x)< b\)で
        \(d\psi(x)\notin\muS(F)\)となる\(C^1\)級関数とする．
        \(\psi\)は\(\supp(F)\)上固有とする．
        このとき，制限射\[
            \RG(M_b;F)\longrightarrow\RG(M_a;F)
        \]は同型である．
    \end{PRP}    
\end{frame}

\begin{frame}
    \frametitle{主張の証明}
    \begin{proof}
        \(\psi\)が固有なので\(
            \Rder{\psi_{!}}F\cong\Rder{\psi_{\ast}}F
        \)である．いま，
        \begin{align*}
            \RG(M_a;F)&\cong\RG({\left]-\infty,a\right[}\mathrel{;}\Rder{\psi_{\ast}}F),\\
            \RG(M_b;F)&\cong\RG({\left]-\infty,b\right[}\mathrel{;}\Rder{\psi_{\ast}}F)        
        \end{align*}
        である．順像の評価より，
        \[
            \muS(\Rder{\psi_{\ast}}F)
            \subset
            \mpb[\psi]\mdiff[\psi]^{-1}\muS(F)
        \]
        なので，仮定の\(d\psi(x)\notin\muS(F)\)から，
        \(a\leqq\psi(x)<b\)に対し，
        \[
            \left(\RG_{[a,b[}\Rder{\psi_{\ast}}F\right)_{\psi(x)}=0
        \]
        となる．
    \end{proof}
\end{frame}


\section*{References}

\begin{frame}[allowframebreaks]{References}
    \begin{thebibliography}{90}\beamertemplatetextbibitems
        \bibitem[GKS12]{GKS12} Stephane Guillermou, 
        Masaki Kashiwara, Pierre Schapira, 
        \textit{Sheaf Quantization of Hamiltonian Isotopies 
        and Applications to Nondisplaceability Problems}, 
        Duke Math. J. 161(2), 201--245, (1 February 2012).
        \bibitem[GS16]{GS16} Stephane Guillermou, Pierre Schapira, 
        \textit{Construction of sheaves on the subanalytic site}, 
        Astérisque, Soc. Math. France, 383 1--60 (2016)
        %\bibitem[Jub13]{Jub13} Benoit Jubin, 
        %\textit{A Microlocal Characterization of Lipschitz Continuity}, 
        %\url{https://doi.org/10.4171/PRIMS/54-4-2}.
        \bibitem[KS90]{KS90} Masaki Kashiwara, Pierre Schapira, 
        \textit{Sheaves on Manifolds}, Springer, 1990.
        \bibitem[Ler45]{Ler45} Jean Leray, 
        \textit{Sur la forme des espaces topologiquest sur
        les points fixes de représentations}, J. Math. Pures Appl. 24 (1945), 95--167.
        \bibitem[SKK73]{SKK73} M. Sato, T. Kawai, M. Kashiwara, 
        \textit{Microfunctions and pseudodifferential equations}, 
        LNM 287, 265--529, Springer-Verlag, 1973.
    \end{thebibliography}
\end{frame}


\begin{comment}
\begin{frame}[noframenumbering]
aaa
\end{frame}
\end{comment}

\end{document}